
% ФОРМУЛЫ
% t_{cp} = K\frac{0,14}{(\frac{I}{I_{cp}})^{0,02}-1}
% t_{cp} = K\frac{13,5}{(\frac{I}{I_{cp}})-1}
% t_{cp} = K\frac{80}{(\frac{I}{I_{cp}})^2-1}
% t_{cp} = K \left( \frac{0,0515}{\left(\frac{I}{I_{cp}}\right)^{0,02} - 1}+0,114 \right)
% t_{cp} = K \left( \frac{19,61}{\left(\frac{I}{I_{cp}}\right)^{2} - 1}+0,491 \right)
% t_{cp} = K \left( \frac{28,2}{\left(\frac{I}{I_{cp}}\right)^{2} - 1}+0,1217 \right)
% t_{cp} =5,8 - 1,35ln \left( \frac{I}{K \cdot I_{cp}} \right)


\KOMAoption{paper}{a4, portrait}
\areaset{180mm}{267mm}
\fancyheadoffset{0pt} % обновляем хедер и футер

\newgeometry{
  includehead,
  includefoot,
  nomarginpar,
  top=6mm,
  bottom=6mm,
  left=22mm,
  right=10mm,
  headheight=14.5pt,
  headsep=3mm,        % Добавить headsep
  footskip=7mm,
}

\color{unimain}{\section[(справочное) Графики обратнозависимых времятоковых характеристик]{(справочное)\\Графики обратнозависимых времятоковых характеристик}\label{app:chars}}
\color{black}

\begin{figure}[h!]
  \centering
  \includegraphics[width=0.5\textwidth,height=0.5\textheight,keepaspectratio]{\apppath/Приложение. Характеристики/img/1.png}
  \caption{Кривая A (Инверсная МЭК)}
  \label{chars:fig1} % Опционально: ссылка на рисунок
\end{figure}

\begin{figure}[h!]
  \centering
  \includegraphics[width=0.5\textwidth,height=0.5\textheight,keepaspectratio]{\apppath/Приложение. Характеристики/img/2.png}
  \caption{Кривая B (Сильно инверсная МЭК)}
  \label{chars:fig2} % Опционально: ссылка на рисунок
\end{figure}

\begin{figure}[h!]
  \centering
  \includegraphics[width=0.55\textwidth,height=0.55\textheight,keepaspectratio]{\apppath/Приложение. Характеристики/img/3.png}
  \caption{Кривая C (Чрезвычайно инверсная МЭК)}
  \label{chars:fig3} % Опционально: ссылка на рисунок
\end{figure}

\begin{figure}[h!]
  \centering
  \includegraphics[width=0.55\textwidth,height=0.55\textheight,keepaspectratio]{\apppath/Приложение. Характеристики/img/4.png}
  \caption{Кривая D (Умеренно инверсная ANSI)}
  \label{chars:fig4} % Опционально: ссылка на рисунок
\end{figure}

\begin{figure}[h!]
  \centering
  \includegraphics[width=0.55\textwidth,height=0.55\textheight,keepaspectratio]{\apppath/Приложение. Характеристики/img/5.png}
  \caption{Кривая E (Сильно инверсная ANSI)}
  \label{chars:fig5} % Опционально: ссылка на рисунок
\end{figure}

\begin{figure}[h!]
  \centering
  \includegraphics[width=0.55\textwidth,height=0.55\textheight,keepaspectratio]{\apppath/Приложение. Характеристики/img/6.png}
  \caption{Кривая F (Чрезвычайно инверсная ANSI)}
  \label{chars:fig6} % Опционально: ссылка на рисунок
\end{figure}


\begin{figure}[h!]
  \centering
  \includegraphics[width=0.85\textwidth,height=0.85\textheight,keepaspectratio]{\apppath/Приложение. Характеристики/img/7.png}
  \caption{Кривая RXIDG}
  \label{chars:fig7} % Опционально: ссылка на рисунок
\end{figure}



%clearpage 
\restoregeometry